%%%%

%%%   Самарский государственный технический университет

%%%   Кафедра прикладной математики и информатики

%%%

%%%   Преамбула для статьи в журнал "Вестник Самарского государственного

%%%   технического университета. Серия: «Физико-математические науки»"

%%%   Ориентирована под MikTEx 2.4 for Win32

%%%

%%%   Хотя сам журнал собирается под GNU/Linux на дистрибутиве TeXLive



\documentclass[11pt,a4paper,twoside]{article}



\usepackage[cp1251]{inputenc}

\usepackage[T2A]{fontenc}

\usepackage[english,russian]{babel}

\usepackage{samgtubib}

\usepackage[dvips]{graphicx}

\usepackage[multiple]{footmisc}



\usepackage{samgtu}

\renewcommand{\VestnikYear}{2009} % Год издания Вестника

\renewcommand{\VestnikNumber}{2\,(19)} % Номер Вестника

\renewcommand{\VestnikMonth}{Ноябрь} % Месяц выхода Вестника

\renewcommand{\VestnikData}{01.11.09} % Дата подписания Вестника в печать

\renewcommand{\VestnikUPL}{26,64} % Усл. печ. л.

\renewcommand{\VestnikUIL}{25,73} % Уч.-изд. л.

\renewcommand{\VestnikRegNumber}{479/09} % Регистрационный номер

\renewcommand{\VestnikOrder}{994} % Номер заказа






%%\usepackage{nccboxes}
%%
%\begin{document}
%%
%
%
%\newcommand{\totop}[1]{\raisebox{6pt}[1pt][2pt]{#1}} 
%\newcommand{\tobot}[1]{\raisebox{-6pt}[1pt][2pt]{#1}} 

\renewcommand{\firstpage}{133}
\renewcommand{\lastpage}{140}


\udk{539.43:621.787} \msc{74G70; 74S05}


\title{Анализ напряж\"eнного состояния в~надрезах полукруглого профиля после опережающего поверхностного пластического деформирования сплошных цилиндрических образцов}
{Анализ напряж\"eннго состояния в~надрезах полукруглого профиля \dots}

\author{М.\,Н.~Саушкин, А.\,Ю.~Куров}{С\,а\,у\,ш\,к\,и\,н~М.\,Н., К\,у\,р\,о\,в~А.\,Ю.}	


\institute{\samgtu.}

\email{E-mails}{msaushkin@gmail.com, alexeykurov@gmail.com}

\otherinf{\fullauthor{Саушкин Михаил Николаевич} \regalia{к.ф.-м.н., доц.} \position{докторант} \dept{каф. прикладной математики и информатики}
\fullauthor{Алексей Юрьевич Куров} \position{студент} \dept{каф. прикладной математики и информатики}
}

\keywords{распределение остаточных напряжений, цилиндрический образец, полукруглый надрез, опережающее упрочнение, метод конечных элементов}

\workabstract{Исследуется распределение остаточных напряжений в~сплошных поверхностно упрочнённых цилиндрических образцах и~образцах с полукруглым надрезом с~использованием метода конечных элементов. 
В качестве исходной информации используется одна и/или две экспериментально определённые компоненты остаточных напряжений в~упрочнённом слое. 
Отмечено, что на перераспределение остаточных напряжений в~области концентратора напряжений существенное влияние оказывает отношение радиуса надреза  к~глубине обнуления осевой компоненты остаточных напряжений. 
}

\engtitle{Analysis of stress state in semicircular profile notches  after preliminary surface plastic deformation of solid cylindrical specimens}


\engauthor{M.\,N.~Saushkin, A.\,Yu.~Kurov}{S\,a\,u\,s\,h\,k\,i\,n~M.\,N., K\,u\,r\,o\,v~A.\,Yu.}


\enginstitute{\engsamgtu.}

\otherenginf{\fullauthor{Mikhail N. Saushkin} \regalia{Ph.\,D. (Phys. {\rm \&} Math.)} \position{Doctoral Candidate}
\dept{Dept. of Applied Mathema\-tics \& Computer Science}
\fullauthor{Alexey Yu. Kurov}
\position{Student}
\dept{Dept. of Applied Mathematics \& Computer Science}}



\engabstract{The distribution of residual stresses of surface hardened solid cylindrical specimens and specimens with a semicircular notch using the finite element method is studied. As the initial information the experimentally determined one and / or two components of residual stresses in a hardened layer are used. It is noted that residual stresses redistribution on the stress concentrator is significantly affected by the ratio of a notch radius to a zero depth of axial component of residual stresses.
}

\engkeywords{residual stresses distribution, cylindrical specimen, semicircular notch, finite element method}

\date{18/I/2012}
\revisiondate{29/II/2012}


\maketitle




Хорошо известно, что зоны концентрации напряжений в~деталях машин и~элементах конструкций являются очагами зарождения усталостных трещин и~хрупких разрушений, развитие которых приводит в~конечном итоге к~их разрушению.
Результаты многих исследователей  свидетельствуют о~положительном влиянии остаточных напряжений на выносливость изделий с~концентраторами напряжений, поэтому для увеличения сопротивления усталости таких изделий в~местах, где может произойти разрушение, при их изготовлении наводят сжимающие остаточные напряжения, как правило, методами поверхностного пластического деформирования. 


%Отметим, что процедуру наведения сжимающих остаточных напряжений в научной и технической литературе принято называть упрочнением.

Упрочнение деталей, имеющих на поверхностях концентраторы напряжений, может производиться после их изготовления, если материал  детали допускает и хорошо воспринимает пластическое деформирование поверхностного слоя на глубину, превышающую соответствующие размеры концентраторов. 
В~противном случае, когда размеры концентратора малы или доступ к нему ограничен,  поле сжимающих остаточных напряжений в~зоне концентратора напряжений стараются наводить при изготовлении.
Иногда на практике изготовлению мелких концентраторов напряжений предшествует упрочнение гладкой детали (опережающее упрочнение поверхности детали).
При таком подходе после упрочнения в~детали возникает неоднородное поле остаточных пластических деформаций и неоднородное по глубине залегания поле остаточных напряжений, а после изготовления концентратора напряжений (фактически "--- удаления части объёма) под действием остаточных пластических деформаций, играющих роль начальных деформаций, в упрочнённом поверхностном слое происходит  перераспределение остаточных напряжений.
Следует отметить, что в~случае цилиндрических деталей метод опережающего упрочнения поверхностей может приводить к большему эффекту, чем непосредственное упрочнение концентратора.

Величина остаточных напряжений в наименьшем сечении цилиндрических деталей с концентраторами непосредственно связана с~сопротивлением многоцикловой усталости  таких изделий, поэтому основной задачей является определение остаточных напряжений в~области концентратора.
Решению этой задачи посвящён цикл работ В.~Ф.~Павлова  с~соавторами [\citen{IvaShaPavl73}--\citen{PavKirChi10} и~др.],
методика которых базируется на  расчёте дополнительных остаточных напряжений в~рамках теории упругости, возникающих в~наименьшем сечении детали после нанесения надреза на упрочнённую поверхность. 
Основные расчётные соотношения (см., например, [\citen{IvaPavlKon92, PavKirIva08}]), которые использовались для определения дополнительных остаточных напряжений цилиндрических образцов диаметром 25~мм и~более, получены методами теории функции комплексного переменного (ТФКП) для осевого сечения детали с~надрезом [\citen{IvaShaPavl73}].
%Для детали с~круговым надрезом они выглядят следующим образом:
%$$
%\sigm_{z\text{д}} (\xi) = \lambda(\xi) + \sum_{k =0, 2, \dots}^{2m} A_k \mu_k(\xi),
%$$
%где $\xi=y$
Для образцов с~диаметром меньше 25~мм использовалась методика, основанная на методе конечных элементов (МКЭ) и~изложенная в~[\citen{Fil85}].
Исходными данными для расчёта дополнительных напряжений в обоих случаях являлись значения осевой компоненты остаточных напряжений гладкой детали в~области сжатия, которые определялись экспериментально.
Суммируя дополнительные остаточные напряжения с напряжениями гладкой детали, авторы работ [\citen{IvaShaPavl73}--\citen{PavKirChi10}] получали остаточные напряжения детали с~надрезом.
Ключевым моментом этой методики является вычисление величины среднеинтегрального осевого остаточного напряжения по толщине упрочнённого слоя, опосредованно зависящей от критической глубины нераспространяющейся трещины.
На основе этой методики установлены зависимости для между остаточными напряжениями и~сопротивлением многоцикловой усталости для ряда деталей с~концентраторами, в~том числе и~для  авиационных резьбовых деталей~[\citen{IvaPavlKon92}].

В работе [\citen{SauKur11}] авторами настоящей работы предложен альтернативный  метод решения этой задачи на основе МКЭ в~вычислительном комплексе {\tt ANSYS}, учитывающий реальное распределение полей остаточных напряжений и~пластических деформаций в~цилиндрическом образце, который подвергся упрочнению.
Исходной информацией при этом является реальное неоднородное по радиусу  распределение остаточных пластических (начальных) деформаций, полученное на основе аналитического решения [\citen{saushkin:11}] по экспериментальной информации о~распределении одной и/или двух компонент остаточных напряжений в упрочнённом слое.
Начальные остаточные пластические деформации в~образце с~надрезом моделируются псевдотемпературными деформациями: полученное поле деформаций для гладкого цилиндрического образца переносится на образец с~надрезом.
После задания начальных деформаций решается задача фиктивной термоупругости относительно неизвестных остаточных напряжений.
Отметим, что распределение полей остаточных напряжений для сплошного цилиндрического образца, полученное МКЭ по псевдотемпературным деформациям,  практически совпадает с~распределением, полученным на основе аналитического решения [\citen{saushkin:11}], что подтверждает корректность решения поставленной задачи.
В~работе~[\citen{saushkin:11}] разработанная методика расчёта была апробирована на образцах из стали 45 и~сплава Д16Т, подвергшихся гидродробеструйной обработке (ГДО, изотропное упрочнение, когда осевая и окружная компонента остаточных напряжений в~области сжатия близки между собой, а~параметр анизотропии упрочнения $\alpha=1$, см. подробнее [\citen{saushkin:11, SauAfaDub08}]).
Проведено сравнение результатов расчётов, выполненных для сплошных цилиндрических образцов и~образцов с~концентраторами, с~экспериментами и расчётами, выполненными по методике~[\citen{PavKirIva08}]. 

В настоящей работе предложенная методика~[\citen{saushkin:11}] применяется как для образцов (рис.~1), предварительное упрочнение которых проводилось методами изотропного упрочнения, так и для образцов с~предварительным анизотропным упрочнением (обкатка роликом, алмазное выглаживание, параметр анизотропии упрочнения \mbox{$\alpha\ne 1$}). При этом исследовалось влияние размера выреза полукруглого профиля на перераспределение остаточных напряжений гладкого цилиндрического образца после предварительного упрочнения поверхности, поэтому в~расчётах в~качестве концентраторов напряжений использовались надрезы с  различным радиусом $\rho$. 

\begin{figure}[b!] \centering \vspace{-3mm}
\includegraphics{./00PIC/obraz.eps}
\caption{Схема образца с круговым надрезом}

\end{figure} 



Типичные картины распределения остаточных напряжений по глубине упрочнённого слоя для сплошных цилиндрических образцов различного радиуса $R$ представлены на рис.~2. 
Здесь значки "--- экспериментальные значения компонент остаточных напряжений~[\citen{PavKirIva08}], на основании которых строится аналитическое решение; сплошные линии "--- значения компонент остаточных напряжений, полученные по авторской методике МКЭ (отметим ещё раз, что они практически совпадают со значениями, полученными по аналитическому решению~[\citen{saushkin:11}]).

\begin{figure}[p!] \centering \small \sl
\includegraphics[scale=0.87]{./00PIC/EI961.eps}

а


\smallskip

\includegraphics[scale=0.87]{./00PIC/30XGSA.eps}

б 


\smallskip

\includegraphics[scale=0.87]{./00PIC/40X.eps}

в

\caption{Распределение остаточных напряжений по глубине $h$ упрочнённого слоя:
\scriptsize
{\sl а\/})~сплав ЭИ 961, $R=5$~мм, алмазное выглаживание, $\alpha=14,6$;
{\sl б\/})~сплав 30ХГСА, $R=7,5$~мм, обкатка роликом, $\alpha=16,6$;
{\sl в\/})~сталь 40Х, $R=12,5$~мм, обкатка роликом, $\alpha=13$;
сплошная линия "--- расчёт; значки "--- экспериментальные данные \cite{PavKirIva08}}

\end{figure} 

\begin{figure}[p!] \centering \small \sl
\includegraphics[scale=0.87]{./00PIC/EI961-conc.eps}

а


\smallskip

\includegraphics[scale=0.87]{./00PIC/30XGSA-conc.eps}

б 


\smallskip

\includegraphics[scale=0.87]{./00PIC/40X-conc.eps}

в

\caption{Распределение остаточных напряжений от дна надреза по толщине слоя $h$ в наименьшем сечении детали:
\scriptsize
{\sl а\/})~сплав ЭИ 961, $R=5$~мм, алмазное выглаживание, $\alpha=14,6$;
{\sl б\/})~сплав 30ХГСА, $R=7,5$~мм, обкатка роликом, $\alpha=16,6$;
{\sl в\/})~сталь 40Х, $R=12,5$~мм, обкатка роликом, $\alpha=13$;
сплошная линия "--- надрез $\rho=0,3$~мм; 
штрихпунктирная линия "--- надрез $\rho=0,5$~мм;
штриховая линия "--- расчёт по методике \cite{IvaShaPavl73, PavKirIva08}}

\end{figure} 

На рис.~3 приведены распределения остаточных напряжений от дна надреза по толщине слоя в наименьшем сечении образца с~концентратором напряжений в виде кругового надреза. Здесь сплошные линии "--- расчёт по авторской методике МКЭ для $\rho=0,3$~мм; 
штрихпунктирные линии "--- расчёт по авторской методике МКЭ для $\rho=0,5$~мм.

Для сопоставимости результатов выполнен сравнительный анализ расчётных значений для напряжения $\sigma_z $ по методике настоящей работы с данными расчёта этой компоненты по приближённой аналитической зависимости, полученной в~\cite{IvaShaPavl73} методами ТФКП для плоской задачи, когда вместо всего цилиндрического образца рассматривалась тонкая пластина "--- осевое сечение цилиндра с~концентратором напряжений.
На рис.~3 данные расчёта величины $\sigma_z $ при $\rho=0,3$~мм  по методике \cite{IvaShaPavl73, PavKirIva08} приведены штриховой линией, при этом наблюдается соответствие данных расчёта по обеим методикам.

Одной из характерных точек эпюры остаточных напряжений является значение на поверхности.
Обозначим  значение компоненты $\sigma_z$ на поверхности цилиндрического образца через $\sigma_{\text{п} z}$, а через
$\sigma_{\text{к} z}$ "--- значение компоненты $\sigma_z$ на поверхности концентратора напряжений.
Аналогично для компоненты $\sigma_\theta$: $\sigma_{\text{п} \theta}$ "--- значение на поверхности образца,
$\sigma_{\text{к} \theta}$ "--- значение на поверхности концентратора.
В~таблице представлены результаты расчётов для 16-ти различных цилиндрических образцов из 6-ти сплавов, подвергшихся различной технологии предварительного упрочнения, с различными концентраторами напряжений (экспериментальные данные брались из [\citen{PavKirIva08}]). 
Здесь $h_0$ "--- глубина  обнуления эпюры $\sigma_z$ для гладкого цилиндрического образца (ещё одна характерная точка).

\begin{table}[b!] \small \centering
\setlength{\tabcolsep}{2pt}
\begin{tabular}{c|c|c|c|c|c|c|c}
\hline
\footnotesize Материал & 
\footnotesize Обработка & 
\footnotesize $R$,  мм & 
\footnotesize $\alpha$ &  
\footnotesize $\rho$, мм & 
\footnotesize $\rho/h0$ & 
\footnotesize $\sigma_{\text{к} z}/ \sigma_{\text{п} z}$ & 
\footnotesize $\sigma_{\text{к} \theta}/ \sigma_{\text{п} \theta}$ \\ \hline
\rule{0mm}{12pt}%
12X18H10T&ГДО&7,50&1,00&0,30&0,73&1,84&0,80 \\
\footnotesize --- // --- &
\footnotesize --- // --- &
\footnotesize --- // --- &
\footnotesize --- // --- &0,50&1,21&0,84&0,32 \\
Д16Т&ГДО&5,00&1,00&0,30&0,64&1,42&0,71 \\
\footnotesize --- // --- &
\footnotesize --- // --- &
\footnotesize --- // --- &
\footnotesize --- // --- &0,50&1,06&0,65&0,26 \\
\footnotesize --- // --- &
\footnotesize --- // --- &7,50&1,00&0,30&0,58&1,42&0,75 \\
\footnotesize --- // --- &
\footnotesize --- // --- &
\footnotesize --- // --- &
\footnotesize --- // --- &0,50&0,97&0,70&0,28 \\
\footnotesize --- // --- &
\footnotesize --- // --- &12,50&1,00&0,30&0,51&1,68&0,80 \\
\footnotesize --- // --- &
\footnotesize --- // --- &
\footnotesize --- // --- &
\footnotesize --- // --- &0,50&0,85&0,81&0,32 \\
\footnotesize --- // --- &
\footnotesize --- // --- &20,00&1,00&0,30&0,47&1,72&0,83 \\
\footnotesize --- // --- &
\footnotesize --- // --- &
\footnotesize --- // --- &
\footnotesize --- // --- &0,50&0,77&0,85&0,34 \\
 Сталь  45 &ГДО&5,00&1,00&0,30&0,64&2,37&1,11 \\
\footnotesize --- // --- &
\footnotesize --- // --- &
\footnotesize --- // --- &
\footnotesize --- // --- &0,50&1,06&0,86&0,34 \\
\footnotesize --- // --- &
\footnotesize --- // --- &7,50&1,00&0,10&0,19&3,36&1,81 \\
\footnotesize --- // --- &
\footnotesize --- // --- &
\footnotesize --- // --- &
\footnotesize --- // --- &0,15&0,29&2,80&1,64 \\
\footnotesize --- // --- &
\footnotesize --- // --- &
\footnotesize --- // --- &
\footnotesize --- // --- &0,30&0,58&1,74&0,86 \\
\footnotesize --- // --- &
\footnotesize --- // --- &
\footnotesize --- // --- &
\footnotesize --- // --- &0,50&0,96&0,94&0,36 \\
\footnotesize --- // --- &
\footnotesize --- // --- &12,50&1,00&0,30&0,55&2,49&1,20 \\
\footnotesize --- // --- &
\footnotesize --- // --- &
\footnotesize --- // --- &
\footnotesize --- // --- &0,50&0,91&1,13&0,45 \\
\footnotesize --- // --- &
\footnotesize --- // --- &25,00&1,00&0,30&0,50&2,44&1,18 \\
\footnotesize --- // --- &
\footnotesize --- // --- &
\footnotesize --- // --- &
\footnotesize --- // --- &0,50&0,83&1,13&0,43 \\
 Сталь  40Х &ГДО&12,50&1,00&0,30&0,68&1,95&0,83 \\
\footnotesize --- // --- & \footnotesize --- // --- & \footnotesize --- // --- & \footnotesize --- // --- &0,50&1,13&0,91&0,33 \\
30XГСА& обкатка  роликом &5,00&35,00&0,30&0,48&3,15&2,43 \\
\footnotesize --- // --- & \footnotesize --- // --- & \footnotesize --- // --- & \footnotesize --- // --- &0,50&0,79&1,68&1,44 \\
\footnotesize --- // --- & \footnotesize --- // --- &7,50&16,60&0,30&0,40&2,94&2,43 \\
\footnotesize --- // --- & \footnotesize --- // --- & \footnotesize --- // --- & \footnotesize --- // --- &0,50&0,66&1,84&1,56 \\
\footnotesize --- // --- & \footnotesize --- // --- &7,50&10,60&0,30&0,67&2,18&1,76 \\
\footnotesize --- // --- & \footnotesize --- // --- & \footnotesize --- // --- & \footnotesize --- // --- &0,50&1,11&0,88&0,76 \\
12X18H10T& обкатка  роликом &5,00&8,20&0,30&0,66&2,34&1,77 \\
\footnotesize --- // --- & \footnotesize --- // --- & \footnotesize --- // --- & \footnotesize --- // --- &0,50&1,10&0,92&0,75 \\
 ЭИ  961 & алмазное выглаживание &5,00&14,60&0,30&0,71&1,14&1,00 \\
\footnotesize --- // --- & \footnotesize --- // --- & \footnotesize --- // --- & \footnotesize --- // --- &0,50&1,19&0,45&0,46 \\
 Сталь  40Х & обкатка  роликом &12,50&13,00&0,30&0,24&3,19&2,58 \\
\footnotesize --- // --- & \footnotesize --- // --- & \footnotesize --- // --- & \footnotesize --- // --- &0,50&0,40&2,60&2,17 \\ \hline
\end{tabular}

\end{table}


На основании представленных   данных можно сделать выводы, которые отчасти повторяют и в то же время дополняют выводы, сформулированные в работе [\citen{SauKur11}].
На перераспределение остаточных напряжений в области концентратора существенное влияние оказывает отношение
 $\rho/h_0$: увеличение $\rho/h_0$ приводит к~уменьшению величины сжимающих остаточных напряжений $\sigma_\theta$~и~$\sigma_z$ на дне концентратора, и~наоборот.
Так, при относительно мелких надрезах ($0,2<\rho/h_0 <0,7$) для исследованных образцов происходит существенное (в~два и более раз) увеличение величины сжимающих остаточных напряжений $\sigma_\theta$ и $\sigma_z$ на дне концентратора по сравнению с~аналогичными напряжениями для образцов без концентраторов напряжений. 
При надрезах, когда $0,75<\rho/h_0 <1,25$,  величины сжимающих остаточных напряжений $\sigma_\theta$~и~$\sigma_z$ на дне концентратора становятся близкими к аналогичным напряжениям в~образце без концентратора. 
При б\'ольших значениях величины $\rho/h_0$ остаточные напряжения на дне концентратора становятся меньше, чем в образце без концентратора. 
Для величины $\sigma_r$ в~образце с  концентратором характерна смена знака,  при этом её значения возрастают (по модулю) на \mbox{1--2}~порядка, а~увеличение $\rho/h_0$ приводит к уменьшению (по модулю) её значений.
Радиусы цилиндрических образцов практически не оказывают влияния на указанную закономерность.




\begin{thebibliography}{9}

\RBibitem{IvaShaPavl73}
\by  Иванов~С.\,И., Шатунов~М.\,П., Павлов~В.\,Ф. 
\paper Определение дополнительных остаточных напряжений в~надрезах на цилиндрических деталях
\inbook Вопросы прочности элементов конструкций
\bookinfo Тр. Куйбышевского авиационного института. Вып.~60
\publaddr Куйбышев
\publ КуАИ
\yr 193
\pages 160--170
\misctransl
\by Ivanov~S.\,I., Shatunov~M.\,P., Pavlov~V.\,F.
\paper Determination of the additional residual stresses in notches on cylindrical parts
\inbook Questions of the Strength of Elements of Aircraft Structures. \rm Issue~60
\publaddr Kuibyshev
\publ KuAI
\yr 193
\pages 160--170

\RBibitem{IvaPavlKon92}
\by  Иванов~С.\,И., Павлов~В.\,Ф., Коновалов~Г.\,В., Минин~Б.\,В.
\book Технологические остаточные напряжения и сопротивление усталости авиационных резьбовых деталей
\publaddr М.
\yr 1992
\totalpages 192
\miscnote Отраслевая библиотека
\misctransl
\by Ivanov~S.\,I., Pavlov~V.\,F., Konovalov~G.\,V., Minin~B.\,V.
\book Technological Residual Stresses and Fatigue Resistance of the Aircraft’s Threaded Components
\publaddr Moscow
\yr 1992
\totalpages 192


\RBibitem{PavBorSur05}
\by Павлов~В.\,Ф, Бордаков~С.\,А., Сургутанова~Ю.\,Н., Каранаева~О.\,В., Хибник~Т.\,А.
\paper Прогнозирование предела выносливости по разрушению деталей, изготовленных методами опережающего поверхностного пластического деформирования
\jour Вестн. Сам. гос. техн. ун-та. Сер. Физ.-мат. науки
\yr 2005
\issue 34
\pages 60--67
\misctransl
\by Pavlov~V.\,F., Bordakov~S.\,A., Surgutanova~Yu.\,N., Karanaeva~O.\,V., Khibnik~T.\,A.
\paper Prediction endurance limits for destruction of details made outstripping methods of surface plastic deformation
\jour Vestn. Samar. Gos. Tekhn. Univ. Ser. Fiz.-Mat. Nauki
\yr 2005
\issue 34
\pages 60--67

\RBibitem{PavKirIva08}
\book Остаточные напряжения и~сопротивление усталости упрочненных деталей с~концентраторами напряжений 
\by Павлов~В.\,Ф., Кирпичев~В.\,А.,  Иванов~В.\,Б.
\publaddr Самара
\publ Сам. науч. центр РАН
\yr 2008
\totalpages 64
\misctransl
\by Pavlov~V.\,F., Kirpichev~V.\,A., Ivanov~V.\,B.
\book Residual Stresses and Fatigue Strength of Hardened Specimens with Stress Concentrators 
\publ Sam. Nauchn. Center RAN.
\publaddr Samara 
\yr 2008
\totalpages 64




\RBibitem{PavKirChi10}
\by Павлов~В.\,Ф., Кирпич\"eв~В.\,А., Чирков~А.\,В., Сем\"eнова~О.\,Ю.
\paper Закономерности распределения дополнительных остаточных напряжений в~упрочнённых цилиндрических деталях с~кольцевыми надрезами полукруглого профиля
\jour Вестн. Сам. гос. техн. ун-та. Сер. Физ.-мат. науки
\yr 2010
\issue 1(20)
\pages 121--126
\misctransl
\by Pavlov~V.\,F., Kirpichev~V.\,A., Chirkov~A.\,V., Semyonova~O.\,Yu.
\paper Regularity of Auxiliary Residual Stresses Distribution into the Surface Treated Cylindrical Details with Annular Semicircular Notches
\jour Vestn. Samar. Gos. Tekhn. Univ. Ser. Fiz.-Mat. Nauki
\yr 2010
\issue 1(20)
\pages 121--126

\RBibitem{Fil85}
\by Филатов~А.\,П.
\preprint Дополнительные остаточные напряжения на дне периодического концентратора, вызванные перераспределением остаточных напряжений гладкой детали 
\publaddr Куйбышев
\yr 1985
\totalpages 24
\miscnote Деп. в~ВИНИТИ 1985. №~5157--в--85. 
\misctransl
\by Philatov~A.\,P.
\yr 1985
\publaddr Kuibyshev
\preprint Additional residual stresses on  periodic stress concentrator bottom  due to  residual stresses on smooth part  redistribution 
\miscnote Deposited at VINITI 1984. No.~5157--v-85
\totalpages 24

\RBibitem{SauKur11}
\by Саушкин~М.\,Н., Куров~А.\,Ю.
\paper Конечно-элементное моделирование распределения остаточных напряжений в~сплошных упроченных цилиндрических образцах и~образцах с~полукруглым надрезом
\jour Вестн. Сам. гос. техн. ун-та. Сер. Физ.-мат. науки
\yr 2011
\issue 3(24)
\pages 72--78
\misctransl
\by Saushkin~M.\,N., Kurov~A.\,Yu.
\paper Finite element modeling of residual stress distribution in solid hardened cylindrical samples and samples with semicircular notch
\jour Vestn. Samar. Gos. Tekhn. Univ. Ser. Fiz.-Mat. Nauki
\yr 2011
\issue  3(24)
\pages 72--78

\RBibitem{saushkin:11}
\by Саушкин~М.\,Н., Радченко~В.\,П., Павлов~В.\,Ф.
\paper Метод расч\"eта полей остаточных напряжений и пластических деформаций в цилиндрических образцах с~уч\"eтом анизотропии процесса поверхностного упрочнения
\jour ПМТФ
\yr 2011
\vol 52
\issue 2
\pages 173--182
\transl
англ. пер.:~
\by Saushkin~M.\,N., Radchenko~V.\,P., Pavlov~V.\,F.
\paper Method of calculating the fields of residual stresses and plastic strains in cylindrical specimens with allowance for surface hardening anisotropy
\jour J. Appl. Mech. Tech. Phys.
\yr 2011
\vol 52
\issue 2
\pages 303--310

\RBibitem{SauAfaDub08}
\by Саушкин~М.\,Н., Афанасьева~О.\,С., Дубовова~Е.\,В., Просвиркина~Е.\,А.
\paper Схема расчёта полей остаточных напряжений в~цилиндрическом образце с~учётом организации процесса поверхностного пластического деформирования
\jour Вестн. Сам. гос. техн. ун-та. Сер. Физ.-мат. науки
\yr 2008
\issue 1(16)
\pages 85--89
\misctransl
\by Saushkin~M.\,N., Afanas'eva~O.\,S., Dubovova~E.\,V., Prosvirkina~E.\,A.
\paper The scheme calculation fields of residual stresses in the cylindrical sample when the process of surface plastic deformation into account taken
\jour Vestn. Samar. Gos. Tekhn. Univ. Ser. Fiz.-Mat. Nauki
\yr 2008
\issue 1(16)
\pages 85--89

\end{thebibliography}






\noindent
\hskip6.5cm \thedate \par\noindent
\hskip6.5cm\therevdate


\makeengtitle


\renewcommand{\baselinestretch}{0.9}

\newpage
%
%\end{document}